\documentclass[11pt,a4paper]{article}

\usepackage[utf8]{inputenc}
\usepackage[T1]{fontenc}
\usepackage[turkish]{babel}
\usepackage{geometry}
\usepackage{array}
\usepackage{booktabs}
\usepackage{longtable}
\usepackage{tabularx}
\usepackage{enumitem}
\usepackage{amsmath}
\usepackage{amssymb}
\usepackage{fancyhdr}
\usepackage{hyperref}
\usepackage{titlesec}
\usepackage{listings}
\usepackage{xcolor}

\geometry{top=22mm,bottom=24mm,left=20mm,right=20mm}
\setlength{\parindent}{0pt}
\setlength{\parskip}{5pt}
\setlist{nosep,leftmargin=16pt}
\renewcommand{\arraystretch}{1.22}

\pagestyle{fancy}
\fancyhf{}
\lhead{\small NEW IDEA EXCHANGE}
\rhead{\small Teknik Rapor}
\cfoot{\small \thepage}
\hypersetup{colorlinks=false,pdfborder={0 0 0}}

\titleformat{\section}{\large\bfseries}{\thesection}{8pt}{}
\titleformat{\subsection}{\normalsize\bfseries}{\thesubsection}{8pt}{}
\titleformat{\subsubsection}{\normalsize\bfseries}{\thesubsubsection}{8pt}{}

\lstset{
  basicstyle=\ttfamily\small,
  columns=fullflexible,
  frame=single,
  breaklines=true,
  tabsize=2,
  showstringspaces=false,
  keywordstyle=\bfseries,
  commentstyle=\itshape
}

\newcommand{\docline}{\noindent\rule{\textwidth}{0.4pt}}
\newcommand{\code}[1]{\texttt{#1}}
\newcommand{\statusok}{\textbf{Uygun}}
\newcommand{\statusrisk}{\textbf{Risk}}
\newcommand{\statusdebt}{\textbf{Teknik borç}}

\begin{document}

\begin{center}
{\Large\bfseries NEW IDEA EXCHANGE}\\[3pt]
{\large\bfseries Sabancı Kurumsal İnovasyon Borsası}\\[6pt]
{\large Teknik Mimari, Mevcut Durum İncelemesi ve Sürdürülebilirlik Raporu}\\[8pt]
\docline\\[6pt]
{\small Tarih: 20 Haziran 2026 \quad Branch: \code{sabanci} \quad Kapsam: Kod incelemesi + teknik rapor}
\end{center}

\tableofcontents
\newpage

\section{Yönetici Özeti}

Bu rapor, \code{C:\textbackslash Users\textbackslash yozid\textbackslash Desktop\textbackslash A\textbackslash App} altında bulunan NEW IDEA EXCHANGE uygulamasının mevcut kod tabanına göre baştan hazırlanmıştır. İnceleme; proje yapısı, frontend mimarisi, state yönetimi, mock veri modeli, AI analiz endpoint'i, Vercel deployment yapısı, kullanıcı akışları, borsa mantığı, yönetici modülleri, dosya görüntüleme katmanı ve teknik borç alanlarını kapsar.

Uygulama, backend'i sınırlı olan fakat ürün demosu açısından geniş kapsamlı bir \textbf{vanilla JavaScript single page application} olarak çalışır. Ana kullanıcı deneyimi \code{public/app.js} içinde tutulur; Sabancı iştirakleri, kullanıcılar, fikirler, ekipler, duyurular, gündem, yarışmalar ve kulüp verileri \code{public/mockData.js} içinde modellenir. Sunucu tarafında iki görev vardır: statik dosya servis etmek ve fikir analizi için \code{/api/ai/analyze-idea} endpoint'ini çalıştırmak.

Ürünün en güçlü tarafı, klasik öneri kutusu yerine fikirleri bir portföy varlığı gibi ele almasıdır. Fikirler fiyat, hacim, değişim, lot, destekçi, AI skoru, ekip bağlantısı, dosya bundle'ı ve ürünleşme seviyesiyle temsil edilir. Yönetici tarafı, yatırım/oy defteri, kategori performansı, karar kuyruğu ve depolama alanıyla desteklenir.

Ana teknik riskler şunlardır:

\begin{enumerate}
  \item \code{public/app.js} çok büyük ve monolitiktir; render, state, event handling, veri dönüşümü ve iş kuralları aynı dosyada birleşmiştir.
  \item State içinde tekrar eden alanlar vardır; örneğin \code{investmentLedger}, \code{ledgerUserFilter}, \code{affiliationFilter} birden fazla kez tanımlanır.
  \item Sabancı branch'inde halen bazı eski \code{İş NEW}, \code{is-new}, \code{Moka} veya eski organizasyon terminolojileri kalmıştır.
  \item Hızlı değerlendirme ekranı görsel olarak \textit{Al/Alma} dilini kullanır; ancak tam borsa Al/Sat defteriyle birebir aynı transaction modeline bağlı değildir.
  \item Gerçek persistence yoktur; çoğu akış tarayıcı belleğinde çalışan demo state ile sınırlıdır.
  \item Otomatik test altyapısı yoktur; mevcut güvence \code{node --check} ve manuel browser smoke test ile sınırlıdır.
\end{enumerate}

Bu risklere rağmen uygulama demo/MVP gösterimi için güçlü bir hikaye sunar. Canlı ürünleşme için önerilen ilk hamle, kodun modülerleştirilmesi, shared data/service katmanı kurulması, state modelinin sadeleştirilmesi ve API mantığının tek kaynaklı hale getirilmesidir.

\section{İnceleme Kapsamı}

İnceleme yerel çalışma ağacı üzerinden yapılmıştır. İncelenen ana dosyalar:

\begin{longtable}{p{0.28\textwidth}p{0.62\textwidth}}
\toprule
\textbf{Dosya} & \textbf{Rol}\\
\midrule
\code{package.json} & Proje metadata'sı ve \code{node server.js} başlangıç script'i.\\
\code{server.js} & Local Node HTTP server, statik dosya servisi, Groq destekli AI analiz endpoint'i.\\
\code{api/ai/analyze-idea.js} & Vercel serverless function karşılığı.\\
\code{vercel.json} & API rewrite ve SPA fallback routing.\\
\code{public/index.html} & SPA shell, fontlar, Lucide, mock data ve app script yükleme sırası.\\
\code{public/mockData.js} & Şirketler, kullanıcılar, fikirler, ekipler, yarışmalar, nav ve başlangıç veri modeli.\\
\code{public/app.js} & Uygulamanın ana state, render, event, işlem ve UI mantığı.\\
\code{public/styles.css} & Tema, responsive grid, borsa, sosyal, stüdyo, ekip ve dosya görüntüleme stilleri.\\
\code{public/assets/*} & Sabancı/iştirak logoları, ülke bayrakları ve görsel varlıklar.\\
\code{reports/*} & Önceki rapor çıktıları ve bu yeni LaTeX rapor kaynağı.\\
\bottomrule
\end{longtable}

Bu rapor kodu değiştirmeyi değil, kodun gerçek halini açıklamayı amaçlar. Rapor içinde belirtilen eksikler, uygulamayı kıran zorunlu hatalar değil; ürünleşme ve bakım kalitesi açısından dikkate alınması gereken teknik borçlardır.

\section{Uygulamanın Ürün Tanımı}

NEW IDEA EXCHANGE, kurum içi fikir, proje, araştırma, şikayet, veri, duyuru, ekip ve ürünleşme süreçlerini tek bir etkileşim alanında toplar. Uygulama Sabancı ekosistemine göre markalanmıştır; iştirakler, ülkeler, şehirler ve kampüsler mock veriyle temsil edilir.

Ürünün ana fikri şudur:

\[
\text{Fikir} \rightarrow \text{Borsa Varlığı} \rightarrow \text{Ekip/Stüdyo} \rightarrow \text{Ürünleşme} \rightarrow \text{Yönetici Kararı}
\]

Bu zincirde fikir yalnızca metin değildir. Her fikir:

\begin{enumerate}
  \item başlık, özet, problem, çözüm ve KPI bilgisi,
  \item şirket/ülke/kampüs bağlantısı,
  \item AI skoru ve stratejik skor,
  \item piyasa fiyatı, değişim, hacim ve lot bilgisi,
  \item yorumlar ve destekçi sayısı,
  \item dosya bundle'ı,
  \item ekip/stüdyo bağlantısı,
  \item ürünleşme aşaması
\end{enumerate}

ile birlikte yaşar.

\section{Genel Mimari}

\subsection{Çalışma Modeli}

Proje iki runtime modeliyle çalışabilir:

\begin{enumerate}
  \item \textbf{Local runtime:} \code{node server.js} ile port \code{4173} üzerinde statik sunum ve AI endpoint.
  \item \textbf{Vercel runtime:} Statik dosyalar + \code{api/ai/analyze-idea.js} serverless function + \code{vercel.json} rewrite.
\end{enumerate}

Basitleştirilmiş akış:

\begin{lstlisting}
Browser
  -> /index.html
  -> /mockData.js
  -> /app.js
  -> render(state)
  -> user action
  -> mutate state
  -> render(state)
\end{lstlisting}

AI analiz akışı:

\begin{lstlisting}
Browser
  -> POST /api/ai/analyze-idea
  -> Groq API varsa model yaniti
  -> Groq yoksa deterministic demo fallback
  -> JSON analysis response
\end{lstlisting}

\subsection{Dosya Bazlı Sorumluluklar}

\begin{longtable}{p{0.24\textwidth}p{0.34\textwidth}p{0.32\textwidth}}
\toprule
\textbf{Katman} & \textbf{Mevcut Sorumluluk} & \textbf{Değerlendirme}\\
\midrule
HTML shell & Root container, fontlar, Lucide, script yükleme. & Basit ve yeterli.\\
CSS & Tema, layout, responsive davranış, kartlar, terminal/borsa UI, dosya modalları. & Geniş fakat parçalı; ileride CSS modülleri önerilir.\\
Mock data & Şirket master data, kullanıcılar, fikirler, ekipler, yarışmalar, nav. & Demo için güçlü; production için API/DB gerekir.\\
App JS & State, render, event delegation, iş kuralları, filtreler, işlem mantığı. & MVP için hızlı; sürdürülebilirlik için bölünmeli.\\
Node server & Static serving + AI endpoint. & Lokal demo için yeterli.\\
Vercel API & AI endpoint'in serverless kopyası. & Mantık server.js ile tekrarlı.\\
\bottomrule
\end{longtable}

\section{Frontend Mimari Analizi}

\subsection{SPA Yaklaşımı}

Uygulama framework kullanmadan tek sayfa uygulama olarak yazılmıştır. Tüm sayfalar string template üreten \code{render*} fonksiyonlarıyla oluşturulur. Kullanıcı etkileşimleri document seviyesinde event delegation ile yakalanır.

Ana render fonksiyonu:

\begin{lstlisting}
function render() {
  document.body.dataset.theme = state.theme;
  if (state.page === "products") {
    state.page = "studio";
    state.studioTab = "products";
  }
  translateAllInState();
  app.innerHTML = state.loggedIn ? renderShell() : renderLogin();
  if (window.lucide) window.lucide.createIcons();
  initAIBubbleDraggability();
}
\end{lstlisting}

Bu yaklaşım hızlı demo üretir; ancak tüm DOM her render'da yeniden yazıldığı için büyük veri setlerinde performans ve focus kaybı riski oluşturur.

\subsection{State Yönetimi}

State global \code{state} nesnesiyle tutulur. Bu nesne login, piyasa, portföy, filtreler, ekipler, kulüpler, gündem, depolama, AI asistan, dosya görüntüleme ve kullanıcı tercihlerini kapsar.

Avantajlar:

\begin{enumerate}
  \item Debug etmek kolaydır.
  \item Demo davranışları hızlıca eklenir.
  \item Backend olmadan çok sayıda ekran tutarlı gösterilebilir.
\end{enumerate}

Riskler:

\begin{enumerate}
  \item State şeması büyüdükçe tutarlılık kontrolü zorlaşır.
  \item Aynı alanların birden fazla tanımlanması son değerin önceki değeri ezmesine neden olur.
  \item Browser refresh sonrası kullanıcı işlemleri kalıcı değildir.
  \item Çok sayıda modül aynı state alanına yazdığı için yan etki takibi zorlaşır.
\end{enumerate}

\section{Veri Modeli}

\subsection{Şirket ve Coğrafya Modeli}

\code{public/mockData.js} içinde \code{affiliationCompanies} listesi Sabancı Holding ve iştiraklerini temsil eder. Her şirket; logo, kısa ad, ülke, şehir, kampüs ve departman alanlarıyla modellenir.

Bu yapı, aşağıdaki filtreleri mümkün kılar:

\begin{enumerate}
  \item ülke portalı,
  \item iştirak filtresi,
  \item şehir/kampüs bazlı duyuru,
  \item ekip ve kişi eşleştirme,
  \item şirket logosu gösterimi.
\end{enumerate}

\subsection{Kullanıcı Modeli}

\code{demoUsers} ve \code{peopleDirectory} iki farklı kullanıcı bağlamı üretir. \code{demoUsers} login ve rol denemesi için kullanılır; \code{peopleDirectory} sosyal, mesajlaşma, ekip üyeliği ve profil kartları için kullanılır.

Kullanıcıda görülen ana alanlar:

\begin{longtable}{p{0.26\textwidth}p{0.58\textwidth}}
\toprule
\textbf{Alan} & \textbf{Amaç}\\
\midrule
\code{id} & Kullanıcı ve ilişki anahtarı.\\
\code{name} & Profil, mesaj ve yatırım defteri gösterimi.\\
\code{companyId} & İştirak ve ülke kapsamı.\\
\code{role} & Profil ve sosyal/ekip görünümü.\\
\code{isManager}, \code{isAdmin} & Yönetici modüllerine erişim.\\
\code{voteCreditBalance} & Destek/oy ve cüzdan davranışı.\\
\code{supportedIdeas} & Kullanıcının desteklediği fikirler.\\
\bottomrule
\end{longtable}

\subsection{Fikir Modeli}

Fikirler \code{initialIdeas} içinde başlar, \code{state.ideas} olarak kopyalanır ve borsa modüllerinde işlenir. Fikir nesnesi hem öneri hem finansal varlık hem ürün adayıdır.

Ana alanlar:

\begin{longtable}{p{0.26\textwidth}p{0.58\textwidth}}
\toprule
\textbf{Alan} & \textbf{Anlam}\\
\midrule
\code{title}, \code{summary} & Kullanıcıya görünen fikir içeriği.\\
\code{problem}, \code{solution} & Karar ve AI analiz bağlamı.\\
\code{companyId}, \code{country} & Portal ve iştirak filtresi.\\
\code{aiScore}, \code{strategicScore} & Önceliklendirme sinyali.\\
\code{marketPrice} & Baz fiyat.\\
\code{marketChange} & Yüzdesel fiyat değişimi.\\
\code{marketVolume} & İşlem/hacim göstergesi.\\
\code{marketTicker} & Borsa sembolü.\\
\code{files} & PDF, PPTX, XLSX, CSV gibi bundle dosyaları.\\
\code{translations} & Çok dilli başlık/özet/problem/çözüm.\\
\bottomrule
\end{longtable}

\section{Borsa ve Fiyatlama Mantığı}

\subsection{Fiyat Formülü}

Kodda \code{marketPrice} fonksiyonu fikir fiyatını baz fiyat ve yüzdesel değişim üzerinden hesaplar:

\begin{lstlisting}
function marketPrice(idea) {
  const base = Number(idea.marketPrice || 100);
  const change = Number(idea.marketChange || 0);
  return Math.max(20, Math.round(base * (1 + change / 100)));
}
\end{lstlisting}

Matematiksel karşılığı:

\[
P_i = \max(20, \operatorname{round}(B_i \cdot (1 + \frac{\Delta_i}{100})))
\]

Burada:

\begin{itemize}
  \item \(P_i\): fikir \(i\) için gösterilen fiyat,
  \item \(B_i\): baz fiyat,
  \item \(\Delta_i\): marketChange yüzdesi.
\end{itemize}

\subsection{Al/Sat İşlemleri}

Borsa ekranında alım ve satım işlemleri \code{marketBudget}, \code{marketHoldings} ve \code{investmentLedger} üzerinde çalışır. Alımda bütçe düşer, holding artar, işlem defterine kayıt düşer. Satımda holding azalır ve bütçe artar.

Alım:

\[
Cash_{t+1} = Cash_t - P_i \cdot q
\]

\[
Holding_{i,t+1} = Holding_{i,t} + q
\]

Satım:

\[
Cash_{t+1} = Cash_t + P_i \cdot q
\]

\[
Holding_{i,t+1} = Holding_{i,t} - q
\]

Mevcut model gerçek emir defteri değildir. Limit emir, piyasa emri, eşleşme motoru, karşı taraf, komisyon ve settlement mantığı demo seviyesinde temsil edilmez.

\subsection{Hızlı Değerlendirme ile Borsa Arasındaki Fark}

Hızlı değerlendirme ekranı kullanıcıya \textit{Sağa Kaydır: Al} ve \textit{Sola Kaydır: Alma} davranışı sunar. Kod seviyesinde bu ekran \code{quickEvalLikes} state'ini ve \code{supportIdea} fonksiyonunu kullanır. Bu, tam borsa alım/satımı değil; oy/destek sinyalidir.

Bu fark ürün dili açısından önemlidir:

\begin{longtable}{p{0.32\textwidth}p{0.28\textwidth}p{0.30\textwidth}}
\toprule
\textbf{Akış} & \textbf{Mevcut Etki} & \textbf{Üretim Önerisi}\\
\midrule
Borsa Al/Sat & Bütçe, holding ve ledger güncellenir. & Transaction-safe ledger.\\
Hızlı Değerlendir Al & Destek/oy sinyali üretir. & İstenirse doğrudan küçük lot alımıyla bağlanmalı.\\
Hızlı Değerlendir Alma & Kartı geçer; negatif işlem defteri üretmez. & Negatif sinyal veya watchlist dışlama olarak modellenmeli.\\
\bottomrule
\end{longtable}

\section{Modül İncelemesi}

\subsection{Giriş ve Yetki}

Giriş ekranı demo exchange key ile başlar. Kullanıcı seçimi rol bazlı menüyü etkiler. Yetki kontrolü \code{canAccess} üzerinden yapılır:

\begin{lstlisting}
function canAccess(item, user = currentUser()) {
  if (item.adminOnly && !user.isAdmin) return false;
  if (item.managerOnly && !user.isManager && !user.isAdmin) return false;
  return true;
}
\end{lstlisting}

Bu model demo için yeterlidir. Üretimde SSO, JWT/session, backend enforcement ve route-level authorization gerekir.

\subsection{Dashboard}

Dashboard; portföy, borsa, hızlı aksiyonlar, veriler ve genel platform sinyallerini birleştirir. Ana amacı kullanıcıyı doğrudan borsa veya yeni başvuru akışına yönlendirmektir.

\subsection{Borsa}

Borsa modülü uygulamanın ana vitrini konumundadır. Fikir listesi, fiyat, değişim, hacim, portföy, liderler, grafik, şirket filtresi ve yeni varlık ekleme akışlarını içerir.

Teknik olarak bu modül şu alanlara bağlıdır:

\begin{enumerate}
  \item \code{state.ideas}
  \item \code{state.marketBudget}
  \item \code{state.marketHoldings}
  \item \code{state.marketDraft}
  \item \code{state.investmentLedger}
  \item \code{filteredMarketIdeas()}
  \item \code{marketPrice()}
\end{enumerate}

\subsection{Veri ve Bilgi}

Veri ve Bilgi modülü veri setleri, kaynaklar ve öneri/şikayet benzeri bildirimleri kapsar. Dosya ekleme ve veri kartları fikir üretim sürecini destekler.

\subsection{Gündem}

Gündem modülü manuel kurumsal gündem akışı olarak tasarlanmıştır. Kod içinde kullanıcıya \textit{AI internete çıkmaz; gündem içerikleri yöneticiler tarafından manuel eklenir} bilgisi verilir. Mevcut ekranda kategori, etiket ve arama filtresi vardır.

Teknik not: Önceki kapsamda beklenen \textit{iç/dış gündem} ve \textit{kurumsal/topluluk} ayrımı bu branch'te görünür bir filtre olarak aktif değildir. Bu raporda bu durum teknik borç olarak işaretlenmiştir.

\subsection{Duyurular}

Duyuru modülü iştirak, ülke, şehir, kampüs ve departman ölçeğinde duyuru yayınlama yaklaşımını destekler. \code{announcementDraft} company/country/campus alanlarıyla çalışır.

\subsection{Sosyal}

Sosyal modül; post, görsel, link, anket, tahminler, leaderboard ve profil etkileşimini içerir. Liderlik tablosu üretim/girişimci, borsacı ve toplam değer mantığına yakın bir demo deneyimi sunar.

\subsection{Stüdyo ve Ekipler}

Stüdyo modülü ürünler ve stüdyolar arasında tab yapısıyla çalışır. \code{products} nav hedefi render aşamasında \code{studio} sayfasına normalize edilir. Ekipler ayrı sayfa olarak korunur ve takım, kulüp, açık pozisyon, ürün bağlantısı ve chat davranışlarını içerir.

\subsection{Yönetici Dashboardu}

Yönetici Dashboardu; toplam oy/lot, toplam değer, en güçlü fikir, aktif kullanıcı, en çok oy alanlar, kategori/stüdyo performansı ve detaylı yatırım defterini gösterir. Veriler \code{state.investmentLedger} ve seed edilmiş demo kayıtlarla birleştirilir.

\subsection{Yönetici Depolama Alanı}

Yönetici Depolama; dosya, not, kaynak, kategori ve tarih alanlarıyla çalışan local-state tabanlı arşivdir. Dosya girişleri gerçek storage'a yüklenmez; dosya adları ve metadata state içinde tutulur.

\subsection{AI Asistan}

AI Asistan widget'ı platform içi veriyle cevap üretir. İnternet araması yapmaz. Trend, gündem, ürün, veri seti ve öneri cevapları \code{handleAIChatResponse} içinde rule-based olarak üretilir. Bu asistan gerçek LLM değil, demo içi yönlendirilmiş yanıt motorudur.

\section{Dosya ve Bundle Yönetimi}

Uygulamada fikir kartlarına ve yönetici depolamaya dosya iliştirme mantığı vardır. Dosya türleri PDF, CSV, XLSX, PPTX, DOC ve görsel önizleme bileşenleriyle temsil edilir.

Kodda görülen önizleme fonksiyonları:

\begin{itemize}
  \item \code{renderPdfPreview}
  \item \code{renderCsvPreview}
  \item \code{renderSheetPreview}
  \item \code{renderDeckPreview}
  \item \code{renderDocPreview}
  \item \code{renderImagePreview}
  \item \code{renderGenericFilePreview}
\end{itemize}

Bu fonksiyonlar gerçek dosya parse etmez; demo seviyesinde dosya adından veya sabit örnek veriden anlamlı önizleme üretir. Üretimde dosya yükleme, storage, virüs tarama, DLP, dosya tipi doğrulama ve erişim kontrolü gerekir.

\section{Backend ve AI Endpoint}

\subsection{Local Server}

\code{server.js} Node built-in \code{http}, \code{fs} ve \code{path} modülleriyle yazılmıştır. Express kullanılmaz. Server iki temel iş yapar:

\begin{enumerate}
  \item public klasöründen statik dosya servis eder,
  \item \code{POST /api/ai/analyze-idea} isteğini işler.
\end{enumerate}

Request body limiti 1 MB olarak kontrol edilir. Path traversal'a karşı \code{sanitizePath} ile public root dışına çıkış engellenir.

\subsection{Vercel Serverless Function}

\code{api/ai/analyze-idea.js}, local server'daki AI analiz mantığının serverless karşılığıdır. Groq API key varsa canlı model çağrısı yapar; yoksa demo-simulation döner.

Desteklenen environment değişkenleri:

\begin{itemize}
  \item \code{GROQ_API_KEY}
  \item \code{GROK_API_KEY}
  \item \code{GROQ_MODEL}
\end{itemize}

Varsayılan model:

\begin{lstlisting}
llama-3.3-70b-versatile
\end{lstlisting}

\subsection{AI Yanıt Sözleşmesi}

Endpoint aşağıdaki şemaya yakın JSON döner:

\begin{lstlisting}
{
  provider,
  generatedAt,
  suggestedTitle,
  score,
  impactScore,
  riskScore,
  feasibilityScore,
  summary,
  strengths,
  missingPoints,
  weaknesses,
  risks,
  costComment,
  pilotSuggestion,
  executiveSummary,
  similarIdeas,
  improvements
}
\end{lstlisting}

Güçlü taraf: API key frontend'e yazılmamıştır. Zayıf taraf: local server ve serverless function aynı mantığı ayrı ayrı taşır. Üretimde \code{analysisService.js} gibi paylaşılan bir modül önerilir.

\section{Vercel Yapılandırması}

\code{vercel.json} iki rewrite içerir:

\begin{lstlisting}
{
  "source": "/api/ai/analyze-idea",
  "destination": "/api/ai/analyze-idea.js"
},
{
  "source": "/(.*)",
  "destination": "/index.html"
}
\end{lstlisting}

Bu yapı SPA fallback için uygundur. Dosya tabanlı statik asset servisinde Vercel'in filesystem önceliği devreye girer; local server'da ise MIME haritası önemlidir.

Teknik not: \code{server.js} içindeki \code{mimeTypes} haritasında \code{.pdf} tanımı yoktur. Local çalışmada PDF asset'leri \code{application/octet-stream} olarak servis edilebilir. Bu kullanıcı deneyimini her zaman bozmaz; ancak doğru header için \code{application/pdf} eklenmelidir.

\section{Güvenlik ve Yetkilendirme Değerlendirmesi}

\subsection{Mevcut Koruma Alanları}

\begin{enumerate}
  \item AI API key frontend'e yazılmaz.
  \item Local server request body büyüklüğünü sınırlar.
  \item Static path normalize edilir ve public root dışına çıkış engellenir.
  \item Manager/admin menüleri frontend rol kontrolüyle gizlenir.
  \item AI fallback, API olmadığı durumda uygulamayı kırmaz.
\end{enumerate}

\subsection{Üretim İçin Eksik Kontroller}

\begin{enumerate}
  \item Backend tarafında gerçek auth enforcement yoktur.
  \item Admin/yönetici state'i client tarafında değiştirilebilir.
  \item Dosya yükleme gerçek storage ve malware scan üzerinden geçmez.
  \item Audit log kalıcı değildir.
  \item Rate limit, abuse prevention ve CSRF/session stratejisi yoktur.
  \item KVKK/veri sınıflandırma kontrolleri UI metni seviyesindedir.
\end{enumerate}

\section{Performans Değerlendirmesi}

\code{public/app.js} yaklaşık 14.7k satırdır. Her render'da \code{app.innerHTML} komple yenilenir. Bu yaklaşım demo ölçekte kabul edilebilir; fakat geniş veri seti, uzun listeler, aktif chat ve dosya önizleme modalında performans sorunları doğurabilir.

Başlıca performans riskleri:

\begin{enumerate}
  \item Global render tüm DOM'u yeniden üretir.
  \item \code{translateAllInState()} her render'da çağrılır.
  \item Büyük mock veri listeleri filtreleme sırasında tekrar tekrar dolaşılır.
  \item Inline style ağırlığı ve dev CSS büyüklüğü bakım maliyetini artırır.
  \item Görseller randomuser ve çeşitli asset kaynaklarından geldiği için network davranışı değişkendir.
\end{enumerate}

\section{Teknik Borç ve Eksik Listesi}

\begin{longtable}{p{0.23\textwidth}p{0.38\textwidth}p{0.29\textwidth}}
\toprule
\textbf{Alan} & \textbf{Bulgu} & \textbf{Öneri}\\
\midrule
State şeması & \code{investmentLedger}, \code{ledgerUserFilter}, \code{affiliationFilter} gibi alanlar birden fazla kez tanımlanıyor. & Tek \code{initialState} şeması ve migration helper.\\
Terminoloji & Sabancı branch'inde bazı \code{İş NEW}, \code{is-new}, eski e-posta domainleri ve eski ürün adları kalmış. & Brand vocabulary taraması ve master text dictionary.\\
Hızlı değerlendirme & UI \textit{Al/Alma}; kod \code{supportIdea} ve \code{quickEvalLikes} ile oy/destek mantığı çalıştırıyor. & Borsa transaction ile ya birleştirilmeli ya da adı \textit{destekle/geç} yapılmalı.\\
Gündem & İç/dış gündem ve kurumsal/topluluk ayrımı bu branch'te görünür filtre değil. & \code{agendaScope}, \code{agendaType} state alanları ve UI filtreleri geri eklenmeli.\\
Products nav & \code{products} clean nav içinde var; render aşamasında \code{studio/products} tabına normalize ediliyor. & Menü dili netleştirilmeli: ayrı sayfa mı, stüdyo tabı mı?\\
API tekrarları & \code{server.js} ve \code{api/ai/analyze-idea.js} benzer analiz mantığını ayrı taşıyor. & Shared service modülü.\\
MIME & Local server \code{.pdf} MIME type tanımıyor. & \code{application/pdf} eklenmeli.\\
Test & Otomatik test suite yok. & Smoke, unit ve basit Playwright akışı.\\
Persistence & Veri tarayıcı belleğinde kalıyor. & DB + API + optimistic UI.\\
Dosya yükleme & Dosyalar gerçek saklama alanına gitmiyor. & Object storage + signed URL + DLP.\\
\bottomrule
\end{longtable}

\section{Önerilen Refactor Planı}

\subsection{Faz 1: Kod Ayrıştırma}

İlk fazda davranış değiştirmeden dosya yapısı sadeleştirilmelidir:

\begin{lstlisting}
public/
  app.js
  state/
    initialState.js
    selectors.js
  modules/
    market.js
    agenda.js
    teams.js
    social.js
    admin.js
    aiAssistant.js
  services/
    analysisClient.js
    fileBundle.js
  ui/
    icons.js
    renderUtils.js
\end{lstlisting}

\subsection{Faz 2: State ve İşlem Defteri}

Borsa işlemleri için tek transaction fonksiyonu önerilir:

\begin{lstlisting}
executeTrade({ ideaId, side, quantity, source }) {
  validateBalance();
  validateHolding();
  mutateCash();
  mutateHolding();
  appendLedger();
  updateMarketSignals();
}
\end{lstlisting}

Hızlı değerlendirme, normal borsa ve detay sayfası aynı fonksiyonu çağırırsa ürün dili ile veri davranışı uyumlu olur.

\subsection{Faz 3: Backend Sınırı}

Üretim ortamında aşağıdaki API sınırları önerilir:

\begin{longtable}{p{0.24\textwidth}p{0.24\textwidth}p{0.42\textwidth}}
\toprule
\textbf{Endpoint} & \textbf{Method} & \textbf{Amaç}\\
\midrule
\code{/api/session} & GET & Kullanıcı, rol ve şirket kapsamı.\\
\code{/api/ideas} & GET/POST & Fikir listeleme ve başvuru.\\
\code{/api/trades} & POST & Al/Sat transaction.\\
\code{/api/agenda} & GET/POST/PATCH & Manuel gündem yönetimi.\\
\code{/api/files} & POST & Signed upload veya metadata kaydı.\\
\code{/api/ai/analyze-idea} & POST & AI analiz.\\
\code{/api/admin/analytics} & GET & Yönetici dashboard metrikleri.\\
\bottomrule
\end{longtable}

\section{Veritabanı Taslak Modeli}

Demo state'i üretim için aşağıdaki tablolara ayrılabilir:

\begin{longtable}{p{0.26\textwidth}p{0.58\textwidth}}
\toprule
\textbf{Tablo} & \textbf{Açıklama}\\
\midrule
\code{companies} & Holding, iştirak, ülke, şehir, kampüs master data.\\
\code{users} & Kullanıcı profili, rol, şirket ve departman.\\
\code{ideas} & Fikir/proje/araştırma/şikayet ana kaydı.\\
\code{idea_scores} & AI, stratejik, topluluk ve risk skorları.\\
\code{market_assets} & Ticker, baz fiyat, değişim, hacim, arz.\\
\code{trade_ledger} & Kullanıcı, fikir, side, quantity, price, timestamp.\\
\code{holdings} & Kullanıcı bazlı portföy.\\
\code{comments} & Fikir, sosyal ve ekip yorumları.\\
\code{teams} & Ürün ekipleri ve rolleri.\\
\code{studios} & İnovasyon stüdyoları.\\
\code{agenda_items} & Manuel kurumsal gündem kayıtları.\\
\code{announcements} & Ölçekli duyurular.\\
\code{admin_storage} & Yönetici kaynakları ve dosya metadata.\\
\code{file_assets} & Dosya adı, tür, storage key, sahiplik ve sınıf.\\
\bottomrule
\end{longtable}

\section{Test Stratejisi}

Mevcut projede otomatik test script'i yoktur. Minimum güvence seti:

\begin{enumerate}
  \item \code{node --check public/app.js}
  \item \code{node --check public/mockData.js}
  \item \code{node --check server.js}
  \item Local smoke: login, borsa, al/sat, gündem, ekip, sosyal, admin storage.
  \item API smoke: \code{POST /api/ai/analyze-idea}.
  \item Vercel smoke: production URL 200, API endpoint 200/400 kontrollü davranış.
\end{enumerate}

Önerilen Playwright senaryoları:

\begin{enumerate}
  \item Demo key ile giriş yap.
  \item Borsa ekranında bir fikir al ve portföy tutarını doğrula.
  \item Aynı fikri sat ve holding azalmasını doğrula.
  \item Global aramada \code{Akbank} ara ve fikir/ekip sonuçlarını gör.
  \item Yönetici kullanıcısıyla dashboard ve depolama alanına eriş.
  \item Normal kullanıcıyla yönetici sayfasına erişemediğini doğrula.
\end{enumerate}

\section{Deployment Notları}

Uygulama son durumda GitHub \code{sabanci} branch'inde tutulmaktadır. Vercel production deploy için mevcut komut akışı:

\begin{lstlisting}
$env:VERCEL_ORG_ID='...'
$env:VERCEL_PROJECT_ID='...'
npx.cmd vercel deploy --prod --yes
\end{lstlisting}

Güvenlik notu: Tokenlar komut satırına veya repo dosyalarına yazılmamalıdır. Vercel token, Groq key ve benzeri bilgiler environment variable veya secret store üzerinden yönetilmelidir.

\section{Önceliklendirilmiş Aksiyon Planı}

\begin{longtable}{p{0.10\textwidth}p{0.30\textwidth}p{0.48\textwidth}}
\toprule
\textbf{Öncelik} & \textbf{Aksiyon} & \textbf{Gerekçe}\\
\midrule
P0 & Brand cleanup & Eski İş NEW/is-new metinleri Sabancı demosunun profesyonel algısını zedeler.\\
P0 & State tekrarlarını temizle & Aynı anahtarların ezilmesi beklenmeyen davranış üretebilir.\\
P1 & Hızlı değerlendirme ile trade engine'i birleştir & Al/Sat dili veri davranışıyla uyumlu olmalı.\\
P1 & Shared AI service & Local ve Vercel endpoint'leri aynı mantığı tek kaynaktan kullanmalı.\\
P1 & MIME ve statik asset kontrolü & PDF ve diğer dosyalar doğru content-type ile servis edilmeli.\\
P2 & Module split & 14.7k satırlık app.js bakım maliyetini düşürmek için ayrıştırılmalı.\\
P2 & Test altyapısı & Deploy öncesi temel akışları otomatik doğrulamak gerekir.\\
P3 & Backend persistence & Demo state üretim verisine dönüştürülmeli.\\
\bottomrule
\end{longtable}

\section{Sonuç}

NEW IDEA EXCHANGE mevcut haliyle güçlü bir demo ürünüdür. Sabancı markasına uyarlanmış iştirak verisi, çok ülkeli kullanıcı seti, fikir borsası, portföy, stüdyo, ekip, sosyal, yarışma, eğitim, etkinlik, yönetici dashboardu ve AI destekli fikir analizi gibi geniş bir ürün hikayesi sunar.

Teknik açıdan en önemli gerçek şudur: uygulama ürün fikrini çok iyi gösterir; fakat kod organizasyonu ve veri kalıcılığı henüz üretim standardında değildir. Raporun önerdiği refactor planı uygulandığında, mevcut demo doğrudan ölçeklenebilir bir kurumsal ürün iskeletine dönüşebilir.

\end{document}
